\begin{abstract}
Although large-scale unsupervised language models (LMs) are capable of acquiring extensive world knowledge and certain reasoning abilities, achieving fine-grained control over their outputs remains challenging due to the inherently unsupervised approach to their training. To improve steerability, current techniques typically rely on human-annotated comparisons of model outputs and subsequently adapt the unsupervised LM to conform to these preferences, most frequently via reinforcement learning from human feedback (RLHF). Nonetheless, RLHF entails a complicated and potentially unstable process, which involves first training a reward model to capture human preferences, and then applying reinforcement learning to update the LM in a way that maximizes the predicted reward while maintaining fidelity to the original model. In this work, we propose a novel parameterization of the reward model used in RLHF, enabling closed-form derivation of the associated optimal policy. This insight allows the conventional RLHF objective to be addressed using a straightforward classification loss. Our proposed method, termed \textit{Direct Preference Optimization} (DPO), exhibits stability, strong empirical performance, and computational efficiency, dispensing with the necessity of language model sampling during fine-tuning and substantially reducing the effort required for hyperparameter tuning. Experimental results indicate that DPO is capable of adapting LMs to human preferences at least as effectively as current leading methods. In particular, DPO surpasses PPO-based RLHF in sentiment control, and delivers comparable or improved results in tasks like summarization and single-turn dialogue, all while offering a much more accessible implementation and training process.
\end{abstract}

\section{Introduction}
Large unsupervised language models (LMs) trained on extensive datasets have demonstrated unexpected competencies~\citep{chowdhery2022palm, brown2020language, touvron2023llama,bubeck2023sparks}. Nonetheless, the human-generated data they rely on encapsulates a broad spectrum of intentions, values, and levels of expertise. Not all these aims and skills are ideal for emulation; for instance, while it is advantageous for an AI coding assistant to \textit{recognize} frequent programming errors to remedy them, when producing code, our preference is for the model to emphasize the (sometimes uncommon) exemplary coding practices embedded within its training set. Analogously, it is beneficial for a language model to be \textit{informed} about a widespread misconception held by half the population, yet undesirable for the model to propagate this error as fact in half of the relevant queries it answers. Thus, the ability to \emph{select and prioritize desired behaviors and responses} from the model’s broad reservoir of \textit{knowledge and skills} is pivotal for constructing AI systems that are reliable, effective, and governable~\citep{ouyang2022training}. Prevailing strategies typically align LMs to human preferences through reinforcement learning (RL); however, as we will demonstrate, the optimization problem faced by RL-based techniques can be precisely solved via a straightforward binary cross-entropy objective, significantly streamlining the preference alignment process.

Broadly, current techniques incorporate targeted behaviors into language models via carefully constructed datasets of human preferences, highlighting actions judged to be safe and beneficial. This alignment phase generally follows large-scale unsupervised pre-training on expansive text corpora. The most direct strategy for preference alignment involves supervised fine-tuning on human-generated demonstrations of desirable outputs, yet the most effective set of approaches comprises reinforcement learning from human or artificial feedback (RLHF/RLAIF; \citep{christiano2017deep,bai2022constitutional}). In RLHF, a reward model is trained on human preference data, after which RL is employed to adjust the language model’s policy to favor responses receiving high reward, while constraining the policy not to diverge too far from the pre-trained model. Although RLHF has yielded language models with remarkable dialogue and programming skills, its process is considerably more intricate than standard supervised learning, requiring multiple LMs to be trained and maintained, as well as in-the-loop sampling from the evolving policy, all of which substantially raise computational requirements.

In this work, we demonstrate a method for directly training a language model to reflect human preferences, bypassing the need for explicit reward modeling and reinforcement learning techniques. We introduce \textit{Direct Preference Optimization (DPO)}, an approach that inherently targets the same goal as established RLHF strategies—namely, reward maximization under a KL-divergence constraint—while offering ease of implementation and transparent training dynamics. Conceptually, the DPO mechanism enhances the relative log-likelihood of preferred outputs compared to less favored alternatives, but crucially, it applies an adaptive, example-specific weighting to mitigate the risk of model collapse observed with simplistic probability ratio objectives. As with prior frameworks, DPO is grounded in a formal preference structure such as the Bradley-Terry model (\cite{bradley1952rankanalysis}), which assesses the congruence between a proposed reward signal and actual human preference judgments. Unlike previous methodologies, which first optimize a reward model using preference-based loss and then align a policy with this reward model, DPO employs a reparameterization to directly express the preference loss in terms of the policy itself. As a result, given access to annotated human preference data for model completions, DPO optimizes the model via a straightforward binary cross entropy loss, thereby recovering the optimal policy with respect to a latent reward shaped by the observed preferences.

The primary contribution of this study is Direct Preference Optimization (DPO), an algorithm that enables the training of language models from human preference data without requiring reinforcement learning. Empirical results indicate that DPO matches or outperforms contemporary techniques, including those based on PPO in RLHF pipelines, on a range of tasks such as sentiment transfer, summarization, and dialogue, when applied to language models of up to 6B parameters.

Large self-supervised language models have demonstrated the ability to solve various tasks in zero-shot settings \citep{radford2019language} or by using few-shot prompts \citep{gpt3,megatron,chowdhery2022palm}. Nevertheless, further performance gains on downstream tasks, along with improved alignment with user objectives, are often realized through fine-tuning these models on collections of human-provided instructions and example outputs \citep{mishra-etal-2022-cross,sanh2022multitask,chung2022scaling,thoppilan2022lamda}. This process, known as `instruction-tuning,' enhances the ability of LLMs to generalize to new instructions outside of those seen during training, thus making them more effective for practical use \citep{chung2022scaling}. Although instruction-tuning is successful, collecting \textit{relative} human preferences between outputs tends to be less labor-intensive than obtaining expert-curated demonstrations. Consequently, later research has focused on fine-tuning LLMs with datasets derived from human preferences, achieving notable improvements in areas like translation \citep{kreutzer-etal-2018-reliability}, summarization \citep{stiennon2022learning,ziegler2020finetuning}, narrative generation \citep{ziegler2020finetuning}, and adherence to instructions \citep{ouyang2022training,ramamurthy2023is}. In these works, the common approach is to first fit a neural reward model to the collected preferences using a model such as the Bradley-Terry model \citep{bradley1952rankanalysis}, and then adapt the language model itself by maximizing this learned reward, typically via reinforcement learning algorithms like REINFORCE \citep{williams1992reinforce}, proximal policy optimization (PPO; \cite{schulman2017proximal}), or related methods \citep{ramamurthy2023is}. Parallel research directions make use of instruction-tuned LLMs and human feedback to generate synthetic preference data, targeting specific attributes including safety and harmlessness \citep{bai2022constitutional}, relying primarily on minimal human supervision provided through textual rubrics used for model annotation. Altogether, these approaches draw from the intersection of two prominent research domains: reinforcement learning-based language model training for various objectives~\citep{Ranzato2015SequenceLT,paulus2018a,wu2018learning}, and the broader field of learning from human preferences \citep{christiano2017deep,kupcsik2018learning}. Despite the advantages of preference-based supervision, the practical difficulties of employing reinforcement learning at scale for large language models persist; this paper introduces a theoretically sound alternative for optimizing over relative preferences that circumvents the need for RL.

Independent of language-based contexts, the problem of learning policies from preferences has been explored in both bandit and reinforcement learning frameworks, with a range of solutions proposed. In contextual bandit scenarios, policy learning via preferences or rankings, as opposed to scalar rewards, is termed the contextual dueling bandit (CDB) problem (\cite{yue2012karmed,dudik2015contextual}). In such settings, where absolute rewards are not available, theoretical work replaces the conventional optimal policy with the notion of a \textit{von Neumann winner}—a policy whose probability of outperforming any competing policy is at least 50\% on expectation \citep{dudik2015contextual}. Notably, in the CDB paradigm, feedback on preferences is typically provided in an online fashion, whereas human preference-based learning generally utilizes a static collection of preference-labeled action pairs collected offline \citep{yan2022human}. Similarly, in \textit{preference-based RL} (PbRL), the learning agent infers policies from binary comparisons induced by an unknown `scoring' function, rather than direct reward signals \citep{BusaFekete2014,ruiz2023dueling}. While multiple PbRL algorithms are available, including those that make use of off-policy preference data, they often proceed by first inferring a latent reward or scoring function, and then separately optimizing the policy with respect to this estimate \citep{jain2013learning,BusaFekete2014,christiano2017deep,sadigh2017active,kupcsik2018learning}. In contrast, we propose an integrated, single-stage policy learning procedure that directly seeks to align the policy with observed preferences.

\section{Preliminaries}\label{section:prelims}

We provide an overview of the RLHF framework introduced by \citeauthor{ziegler2020finetuning} and further extended in \citep{stiennon2022learning, bai2022training, ouyang2022training}. The standard RLHF methodology is typically structured into three primary steps: (1) supervised fine-tuning (SFT); (2) collection of preference data and subsequent reward model training; and (3) reinforcement learning for policy improvement.

\textbf{SFT}: The RLHF process customarily initiates with supervised fine-tuning of a pre-trained language model using curated, high-quality datasets that reflect the target tasks (e.g., conversation, summarization), resulting in the SFT policy $\pi^\text{SFT}$.

\textbf{Reward Modelling Phase}: During the reward modeling phase, the SFT model is conditioned on input prompts $x$ to generate response pairs $(y_1, y_2)\sim \pi^\text{SFT}(y \mid x)$. Human annotators are then asked to select their preferred response from each pair, indicating a preference $y_w\succ y_l \mid x$ where $y_w$ is the favored completion and $y_l$ is the less preferred, among $(y_1, y_2)$. These preferences are posited to arise from an unobserved reward function $r^*(y, x)$, to which we have no direct access. Several preference modeling strategies have been employed, with the Bradley-Terry (BT) \cite{bradley1952rankanalysis} model being especially prominent, though more general ranking schemes such as the Plackett-Luce model \citep{plackett1975analysis, luce2012individual} are compatible, particularly when multiple ranked responses are provided. The BT model characterizes the human preference probability distribution $p^*$ as follows:

\begin{equation}\label{eq:bradley-terry}
    p^*(y_1\succ y_2 \mid x)=\frac{\exp\left(r^*(x, y_1)\right)}{\exp\left(r^*(x, y_1)\right) + \exp\left(r^*(x, y_2)\right)}.
\end{equation}

Suppose we are provided with a fixed dataset of comparison samples $\mathcal{D}=\bigl\{x^{(i)}, y_w^{(i)}, y_l^{(i)}\bigr\}_{i=1}^N$, which are drawn according to $p^*$. We can then introduce a parameterized reward function $r_{\phi}(x, y)$ and fit its parameters by maximizing the likelihood. Treating the task as binary classification, the negative log-likelihood objective can be written as follows:

\begin{equation}\label{eq:reward_model}
    \mathcal{L}_R(r_{\phi}, \mathcal{D}) = -\mathbb{E}_{(x, y_w, y_l)\sim \mathcal{D}}\bigl[\log \sigma(r_{\phi}(x, y_w)- r_{\phi}(x, y_l))\bigr]
\end{equation}

where the logistic function is denoted by $\sigma$. When working with LMs, $r_{\phi}(x, y)$ is generally initialized using the SFT model $\pi^\text{SFT}(y \mid x)$, and is augmented with an additional linear layer above the final transformer layer to yield a single scalar as the reward prediction~\cite{ziegler2020finetuning}. To control reward variance, it is customary to normalize $r_\phi$, ensuring $\mathbb{E}_{x,y\sim \mathcal{D}}\left[r_\phi(x, y)\right] = 0$ for all $x$.

\textbf{RL Fine-Tuning Phase}: In this phase, feedback from the learned reward function is employed to adjust the language model. Building on previous work~\citep{jaques2017sequence, jaques2020human}, the objective is framed as

\begin{equation}\label{eq:RL}
\max_{\pi_{\theta}}  \mathbb{E}_{x\sim \mathcal{D}, y\sim \pi_{\theta}(y \mid x)}\bigl[r_{\phi}(x, y)\bigr] - \beta\mathbb{D}_{\textrm{KL}}\bigl[\pi_{\theta}(y\mid x)\mid \mid \pi_\text{ref}(y\mid x)\bigr],
\end{equation}

where $\beta$ regulates how much the learned policy $\pi_\theta$ is allowed to diverge from a reference policy $\pi_\text{ref}$ (typically the original SFT policy $\pi^\text{SFT}$). In practice, the starting point for $\pi_\theta$ is also $\pi^\text{SFT}$. This constraint is vital to ensure the model remains close to the distribution where the reward function is reliable and preserves response diversity, thereby avoiding collapse to a small set of high-reward outputs. Since language generation involves discrete outputs, direct differentiation through this objective is not feasible, so reinforcement learning is typically employed for optimization. The common procedure~\citep{ziegler2020finetuning, stiennon2022learning, bai2022training, ouyang2022training} is to define the reward as ${r(x, y) = r_{\phi}(x, y) -\beta (\log \pi_{\theta}(y\mid x) - \log \pi_\text{ref}(y\mid x))}$ and then perform optimization using PPO~\cite{schulman2017proximal}.

\section{Direct Preference Optimization}\label{sec:DPO}

Given the limitations inherent to reinforcement learning methods—especially when scaling to large models for language fine-tuning—we aim to formulate a more streamlined, preference-driven policy optimization approach. Departing from traditional RLHF, which first estimates a reward function and then applies reinforcement learning, our method exploits a specific reward parameterization that admits a closed-form solution for the optimal policy, thus eliminating the need for an RL optimization cycle. In the following sections, we detail our principal insight: by making use of a direct analytical relationship between reward functions and their associated optimal policies, we can re-express a reward-based loss as a policy-based loss. This change of variables obviates the need to explicitly fit a reward model, yet remains compatible with established human preference modeling frameworks such as the Bradley-Terry model. The resulting

\textbf{Deriving the DPO objective.} We begin our derivation from the standard RL objective as formulated in Eq.~\ref{eq:RL}, involving a generic reward function $r$. Building on earlier research~\citep{peters2007reinforcement, peng2019advantage, korbak2022reinforcement, go2023aligning}, one can demonstrate that the solution maximizing the expected reward, subject to a KL constraint as in Eq.~\ref{eq:RL}, is given by:

\begin{equation}\label{eq:op_policy}
    \pi_r(y\mid x) = \frac{1}{Z(x)}\pi_\text{ref}(y\mid x)\exp\left(\frac{1}{\beta}r(x, y)\right),
\end{equation}

where the normalization constant $Z(x) =\sum_{y}\pi_\text{ref}(y\mid x)\exp\left(\frac{1}{\beta}r(x, y)\right)$ is known as the partition function. The details of this derivation are provided in Appendix \ref{app:derivation1}. Notably, even if we substitute the MLE-based estimate $r_{\phi}$ for the actual reward $r^*$, estimating $Z(x)$ remains computationally intensive~\citep{korbak2022reinforcement, go2023aligning}, limiting the practicality of this policy form. Nevertheless, Eq.~\ref{eq:op_policy} can be manipulated to isolate $r(x,y)$ in terms of $\pi_r$, $\pi_\text{ref}$, and the yet-to-be-determined partition function $Z(\cdot)$. By applying the logarithm to both sides of Eq.~\ref{eq:op_policy} and rearranging, we obtain:

\begin{equation}\label{eq:main_eq}
    r(x,y) =\beta \log \frac{\pi_r(y\mid x)}{\pi_\text{ref}(y\mid x)} + \beta \log Z(x).
\end{equation}

This reparameterization holds for the true reward $r^*$ and its associated optimal policy $\pi^*$. Importantly, the Bradley-Terry framework relies solely on reward differences for two outcomes; that is, the probability of preference is defined as ${p^*(y_1 \succ y_2 \mid x) = \sigma(r^*(x, y_1) - r^*(x, y_2))}$. When substituting Eq.~\ref{eq:main_eq} into the Bradley-Terry preference formula (Eq.~\ref{eq:bradley-terry}), the normalization terms $Z(x)$ are eliminated, allowing the human preference likelihood to be described entirely by $\pi^*$ and $\pi_\text{ref}$. Therefore, the optimal RLHF policy $\pi^*$ under the Bradley-Terry assumption fulfills:

\begin{equation}\label{eq:objective}
    p^*(y_1\succ y_2 \mid x)=\frac{1}{1 + \exp\left(\beta \log \frac{\pi^*(y_2\mid x)}{\pi_\text{ref}(y_2\mid x)} - \beta \log \frac{\pi^*(y_1\mid x)}{\pi_\text{ref}(y_1\mid x)}\right)}
\end{equation}

A complete derivation is presented in Appendix~\ref{app:derivation2}. While Eq.~\ref{eq:objective} employs the Bradley-Terry model, the approach can be generalized to alternative ranking models such as Plackett-Luce~\citep{plackett1975analysis, luce2012individual}, as elaborated in Appendix~\ref{app:plackett_luce_models}.

With human preference probabilities now characterized directly via the policy rather than a reward function, it becomes possible to construct a likelihood maximization objective with respect to a parametrized policy $\pi_\theta$. This objective parallels the classical reward modeling loss (see Eq.~\ref{eq:reward_model}), resulting in:

\begin{equation}\label{eq:optimum_model}
    \mathcal{L}_\text{DPO}(\pi_{\theta}; \pi_\text{ref}) = -\mathbb{E}_{(x, y_w, y_l)\sim \mathcal{D}}\left[\log \sigma \left(\beta \log \frac{\pi_{\theta}(y_w\mid x)}{\pi_\text{ref}(y_w\mid x)} - \beta \log \frac{\pi_{\theta}(y_l\mid x)}{\pi_\text

In this formulation, we estimate an implicit reward function by employing a distinct parameterization, where the resulting optimal policy corresponds directly to $\pi_\theta$. Additionally, our approach, being mathematically equivalent to fitting a reparametrized Bradley-Terry model, inherits specific theoretical guarantees—including consistency under appropriate assumptions about the distribution of preference data \cite{bong2022generalized}. We provide further analysis of these theoretical characteristics of DPO and comparisons with related methodologies in Section~\ref{sec:theory}.

\textbf{What is the mechanism behind the DPO update?} To obtain mechanistic insight into DPO, we can examine the gradient of its loss, denoted by $\mathcal{L}_\text{DPO}$. The gradient with respect to parameters $\theta$ takes the following form:

\begin{multline*}\label{eq:gradient}
    \nabla_\theta \mathcal{L}_\text{DPO}(\pi_\theta;\pi_\text{ref}) = \\ -\beta\mathbb{E}_{(x, y_w, y_l) \sim \mathcal{D}} \bigg[\underbrace{\sigma(\hat{r}_\theta(x, y_l) - \hat{r}_\theta (x, y_w))}_\text{greater weight for misranked rewards}\bigg[\underbrace{\nabla_\theta\log \pi(y_w \mid x)}_\text{promote $y_w$'s likelihood} - \underbrace{\nabla_\theta\log\pi(y_l \mid x)}_\text{penalize $y_l$'s likelihood}\bigg]\bigg],
\end{multline*}

Here, the implicit reward $\hat{r}_\theta(x, y)$ is given by $\beta \log \frac{\pi_\theta(y \mid x)}{\pi_\text{ref}(y \mid x)}$, as defined by both the language model $\pi_\theta$ and the reference policy $\pi_\text{ref}$ (see Section~\ref{sec:theory} for further discussion). The gradient of $\mathcal{L}_\text{DPO}$ thus serves to boost the probability of preferred responses $y_w$, while reducing the probability of less-preferred responses $y_l$. Crucially, each training instance is assigned a weight reflecting the degree to which the implicit reward model $\hat{r}_\theta$ incorrectly ranks $y_l$ above $y_w$, modulated by $\beta$ to account for the influence of the KL regularization. Empirical results demonstrate that this weighting scheme is essential: removing it—as in a simple version of this approach—can result in pathological behavior by the language model (refer to Appendix Table~\ref{tab:unlikelihood_generations}).

\textbf{DPO summary.} The DPO workflow generally involves: 1) Drawing sample completions $y_1, y_2 \sim \pi_\text{ref}(\cdot \mid x)$ for each prompt $x$, using human judgments to annotate these and build the offline preference dataset $\mathcal{D} = \{x^{(i)}, y_w^{(i)}, y_l^{(i)}\}_{i=1}^N$, and 2) updating the language model $\pi_\theta$ by minimizing $\mathcal{L}_\text{DPO}$ with respect to the specified $\pi_\text{ref}$, $\mathcal{D}$, and target $\beta$. In applied scenarios, rather than generating new samples and obtaining fresh human preferences, one typically leverages publicly available preference datasets. Since these datasets are usually generated with $\pi^\text{SFT}$, we set $\pi_\text{ref} = \pi^\text{SFT}$ whenever this is possible. In cases where $\pi^\text{SFT}$ is absent, $\pi_\text{ref}$ is initialized by maximizing the likelihood over the preferred completions ${(x, y_w)}$; concretely, this means ${\pi_\text{ref} = \argmax_{\pi}\mathbb{E}_{x, y_w \sim \mathcal{D}}\left[\log \pi(y_w \mid x)\right]}$. This step is designed to address distributional mismatches between the intractable true reference distribution and the $\pi_\text{ref}$ actually used in DPO. Additional implementation specifics and choices of hyperparameters are available in Appendix~\ref{app:implementation}.

\section{Theoretical Analysis of DPO} Here, we further interpret the DPO approach, establish theoretical foundations, and highlight how DPO remedies limitations commonly found in actor-critic RLHF techniques like PPO~\cite{schulman2017proximal}.

\label{sec:theory}

\subsection{Your Language Model Is Secretly a Reward Model} DPO manages to eliminate both explicit reward modeling and separate RL-based policy training by consolidating these into a single maximum likelihood framework. Specifically, the objective in Eq.~\ref{eq:main_eq} mirrors a Bradley-Terry model, with rewards defined as $r^*(x, y) = \beta \log\frac{\pi^*_\theta(y \mid x)}{\pi_\text{ref}(y \mid x)}$, and training our parametric policy $\pi_\theta$ becomes analogous to optimizing a reward model as in Eq.~\ref{eq:reward_model} via a variable transformation. This section will elaborate the theoretical underpinnings of this reparameterization, demonstrating that it introduces no restrictions on the expressiveness of the learned reward models and supports perfect recovery of the optimal policy. To proceed, we introduce a formal equivalence notion among reward functions.

\begin{definition}
Two reward functions $r(x, y)$ and $r'(x, y)$ are said to be equivalent if and only if there exists some function $f$ such that ${r(x, y)-r'(x, y) = f(x)}$.
\end{definition}

This equivalence relation is readily verified and naturally divides the collection of reward functions into equivalence classes. Based on this, we can present the following lemmas:

\begin{lemma}\label{lemma:same_prefrence}
Within the Plackett-Luce preference model, which includes the Bradley-Terry model as a special case, any two reward functions that belong to the same equivalence class yield identical preference distributions.
\end{lemma}

\begin{lemma}\label{lemma:same_policy}
    For the constrained RL setting, two reward functions drawn from the same equivalence class result in the same optimal policy.
\end{lemma}

Proofs for these results are straightforward and are provided in Appendix \ref{app:lemma1}. The first lemma highlights a well-established issue of under-identification associated with the Plackett-Luce family of preference models \cite{plackett1975analysis}. This under-determination necessitates incorporating suitable identifiability constraints in order to obtain well-defined maximum likelihood estimates from Eq. \ref{eq:reward_model} \cite{bong2022generalized}. The second lemma demonstrates that all reward functions from a particular equivalence class will prescribe the same optimal policy, so, for practical purposes, recovering any representative reward function from the optimal class suffices. We establish the following Theorem, with proof in Appendix~\ref{app:thm1}:

\begin{theorem}\label{thm:main}
    Provided certain mild conditions hold, every reward equivalence class compatible with Plackett-Luce (including, specifically, the Bradley-Terry model) admits a representation of the form ${r(x, y) = \beta \log \frac{\pi(y\mid x)}{\pi_\text{ref}(y\mid x)}}$ for an appropriate choice of model $\pi(y\mid x)$ and a fixed reference model $\pi_\text{ref}(y \mid x)$.
\end{theorem}

\begin{sproof}
    Let us consider any given reward function $r(x, y)$, which induces an optimal model $\pi_r(y \mid x)$ as described by Eq. \ref{eq:op_policy}. We intend to demonstrate that some member of $r$'s equivalence class can be expressed via the proposed reparameterization. We introduce the projection function $f$ as follows:
\begin{equation}
    f(r; \pi_\text{ref}, \beta)(x, y) = r(x, y) - \beta\log\sum_{y}\pi_\text{ref}(y\mid x)\exp\left(\frac{1}{\beta}r(x, y)\right)
\end{equation}
This operation $f$ normalizes the reward function by subtracting the log of the associated partition function under $\pi_r$. Since this normalization term depends solely on $x$, the resulting $f(r; \pi_\text{ref}, \beta)(x, y)$ resides in the same equivalence class as $r(x, y)$. Substituting $r$ with the right-hand side of Eq.~\ref{eq:main_eq} (which is valid for any reward function), we observe that $f(r; \pi_\text{ref}, \beta)(x, y) = \beta \log \frac{\pi_r(y\mid x)}{\pi_\text{ref}(y\mid x)}$. Thus, the projection $f$ selects a representative in $r$'s equivalence class with the required structure, implying that the reparameterized reward model does not sacrifice generality.
\end{sproof}

Theorem~\ref{thm:main} can also be interpreted as identifying the specific reward function, from each equivalence class, chosen by the DPO reparameterization—namely, the reward function that fulfills

\begin{equation}\label{eq:lag_p}
     \sum_{y}\underbrace{\pi_\text{ref}(y\mid x)\exp\left(\frac{1}{\beta}r(x, y)\right)}_{=\pi(y\mid x)\text{, using Thm.~\ref{thm:main} reparam.}} = 1,
\end{equation}

meaning that $\pi(y\mid x)$ is indeed a valid probability distribution, being positive and normalized to sum to unity. Examining Eq.~\ref{eq:op_policy}, it becomes apparent that Eq.~\ref{eq:lag_p} functions as the partition function associated with the optimal policy determined by the reward function $r(x, y)$. The central insight underpinning the DPO algorithm is that, by enforcing suitable constraints on the inherently under-specified Plackett-Luce family (including the Bradley-Terry model) of preference representations, one can maintain the expressivity of possible reward models while ensuring that the optimal policy described by Eq.~\ref{eq:op_policy} remains analytically manageable for any prompt $x$.

\subsection{Instability of Actor-Critic Algorithms} The presented framework further allows us to pinpoint sources of instability in standard actor-critic approaches utilized in RLHF pipelines, such as PPO. Focusing on the RL fine-tuning stage as described in Section \ref{section:prelims}, we can relate this process to the control-as-inference perspective on constrained RL described in \ref{eq:RL}, following \cite{levine2018reinforcement}. Assuming a parameterized policy $\pi_{\theta}(y\mid x)$, the objective is to minimize $\mathbb{D}_{\text{KL}}[\pi_{\theta}(y|x) \mid \mid \pi^*(y\mid x)]$, where $\pi^*$ denotes the optimal policy, derived from Eq.~\ref{eq:optimum_model} and associated with the reward function $r_{\phi}(y, x)$. Some algebraic manipulations reduce this to the following optimization problem:

\begin{equation}\label{eq:AC}
    \max_{\pi_{\theta}}\mathbb{E}_{\pi_{\theta}(y\mid x)}\bigg[\underbrace{r_{\phi}(x, y) -\beta\log\sum_{y}\pi_\text{ref}(y\mid x)\exp\left(\frac{1}{\beta}r_{\phi}(x, y)\right)}_{f(r_{\phi}, \pi_\text{ref}, \beta)} - \underbrace{\beta\log\frac{\pi_{\theta}(y\mid x)}{\pi_\text{ref}(y\mid x)}}_{\text{KL}}\bigg]
\end{equation}

The objective considered here is identical to that targeted in earlier studies 
\citep{ziegler2020finetuning, stiennon2022learning, bai2022training, ouyang2022training}, in which the DPO-equivalent reward structure is used for the reward function class $r_{\phi}$. In this formulation, the normalization component of $f(r_{\phi}, \pi_\text{ref}, \beta)$ can be viewed as representing the soft value function associated with the reference policy $\pi_\text{ref}$. Although this normalization does not alter the location of the optimum, its absence may result in policy gradient estimates with significant variance, thereby potentially destabilizing training. One approach to address normalization involves employing a learned value function, though this method can pose substantial optimization challenges. Another common strategy in previous literature is to normalize the reward via a human completion baseline, which effectively serves as a single-sample Monte Carlo estimate for the normalization. In contrast, the DPO reparameterization constructs a reward function that is intrinsically baseline-free.

\section{Experiments}
This section presents empirical studies assessing the effectiveness of DPO for training policies directly from preference data. Initially, we explore a controlled text generation environment to investigate DPO’s efficiency in balancing the dual objectives of reward maximization and KL-divergence minimization relative to the reference policy, in comparison with established preference-based algorithms such as PPO. Subsequently, we analyze DPO’s capacity on larger-scale models and more challenging RLHF benchmarks, which encompass tasks like summarization and dialogue generation. The experiments indicate that, even with minimal hyperparameter adjustment, DPO generally matches or outperforms strong baselines such as PPO-based RLHF and approaches relying on selecting the best of $N$ sampled trajectories under a learned reward. Preceding the presentation of empirical results, we detail the experimental protocol; further experimental specifics can be found in Appendix~\ref{app:exp_details}.

\textbf{Tasks.} We investigate three distinct tasks centered on open-ended text generation. In each case, the various algorithms aim to optimize a policy based on a dataset of preferences, denoted $\mathcal{D}=\bigl\{x^{(i)}, y_w^{(i)}, y_l^{(i)}\bigr\}_{i=1}^N$. For \textbf{controlled sentiment generation}, the input $x$ represents the opening of a movie review drawn from the IMDb dataset \cite{maas-EtAl:2011:ACL-HLT2011}, and the objective for the policy is to generate a continuation $y$ exhibiting positive sentiment. To facilitate a controlled evaluation, we construct preference pairs by generating outputs and ranking them according to a pre-trained sentiment classifier, ensuring that $p(\text{positive}\mid x,y_w)>p(\text{positive}\mid x,y_l)$. In the case of SFT, we further train GPT-2-large on the training set of IMDB reviews until the model converges; more implementation details can be found in App~\ref{app:sentiment_details}. In the \textbf{summarization} task, $x$ is a Reddit forum post, and the goal is to produce a summary $y$ encapsulating the main arguments or themes. Following established methodologies, we employ the Reddit TL;DR summarization dataset \citep{volske-etal-2017-tl}, complemented by human preference data from \citeauthor{stiennon2022learning}. The SFT baseline is a model adapted through fine-tuning on summaries authored by humans\footnote{\url{https://huggingface.co/CarperAI/openai_summarize_tldr_sft}} and further optimized with the TRLX \citep{leandro_von_werra_2023_7790115} toolkit within an RLHF paradigm. The human preference data utilized for comparison was originally collected by \citeauthor{stiennon2022learning} from outputs of a different, yet comparably fine-tuned, SFT model. Lastly, in the context of \textbf{single-turn dialogue}, each $x$ consists of a user inquiry, ranging in topic from scientific queries (e.g., astrophysics) to personal advice. The policy's challenge is to generate a relevant and useful response $y$; we rely on the Anthropic Helpful and Harmless dialogue dataset \citep{bai2022training}, which contains approximately 170,000 dialogues between humans and an automated assistant. Each conversation concludes with two responses generated by a large language model, alongside human preference annotations indicating the favored reply. Unlike the previous tasks, no pre-existing SFT model is available here, so we develop the SFT model by fine-tuning a generic language model on the subset of preferred responses.

\textbf{Evaluation.} Our experimental framework employs two distinct evaluation strategies. To measure how well each algorithm achieves the constrained reward maximization goal in a controlled sentiment generation scenario, we chart the set of reward-KL-divergence pairs that each method can attain. This analysis is feasible since we possess the exact reward function—namely, a sentiment classifier—which allows us to compute the optimal performance trade-off curve. Conversely, when the true reward is unavailable, as is common in practical settings, we assess each approach by its \textit{win rate} in head-to-head comparisons against a baseline, using GPT-4 to stand in for human judges regarding summary quality (for summarization) and helpfulness (for single-turn dialogue). For the summarization task, we adopt the reference summaries from the test dataset as baselines; for dialogue, we compare to the annotated preferred responses. While prior research indicates that large language models can provide more consistent automated evaluations than traditional metrics \citep{Chen2023ExploringTU}, we also include a human assessment (see Sec.~\ref{sec:human-judgments}) to validate our reliance on GPT-4-based evaluation. Our findings indicate that GPT-4’s evaluations are highly consistent with human raters; in fact, the level of concordance between humans and GPT-4 matches or exceeds the agreement among different human annotators.

\textbf{Methods.} Alongside DPO, our comparison covers a range of established techniques for preference-based language model training. For baseline prompting strategies, we utilize zero-shot prompting with \textbf{GPT-J} \citep{gpt-j} in the summarization context and 2-shot prompting with \textbf{Pythia-2.8B} \citep{biderman2023pythia} for dialogue. Further, we assess both the \textbf{SFT} model and the \textbf{Preferred-FT} approach, where the model is further tuned via supervised learning to produce the selected completion $y_w$—originating from either SFT outputs (for sentiment and summarization) or from a generic LM (for dialogue). We also investigate the \textbf{Unlikelihood}~\citep{welleck2019neural} strategy, which aims to increase the likelihood of $y_w$ while explicitly reducing the probability of $y_l$; a weighting parameter $\alpha\in[0,1]$ is optionally applied to the penalization term. Additionally, we test \textbf{PPO} \citep{schulman2017proximal} based on a reward model trained from preference annotations, as well as an oracle version, \textbf{PPO-GT}, which directly uses the true reward signal in the controlled sentiment environment. For our sentiment-based experiments, we employ two PPO-GT variants: a standard implementation \cite{leandro_von_werra_2023_7790115} and a customized version that introduces reward normalization and optimized hyperparameter tuning to enhance results (the same modifications are applied to PPO with learned rewards). Lastly, we include a \textbf{Best of $N$} benchmark, which involves generating $N$ candidates from the SFT model (or from Preferred-FT in the dialogue case) and selecting the response with the highest predicted score as per the learned reward function. Although this strategy can yield superior outcomes by separating reward estimation from PPO dynamics, it quickly becomes infeasible in practice for modest $N$ values since it requires generating and ranking $N$ completions per test instance.

\subsection{How well can DPO optimize the RLHF objective?}

In standard RLHF frameworks, the reward maximization objective subject to a KL constraint seeks to encourage policies that maximize the reward without straying significantly from the reference policy. Thus, to make fair comparisons among different algorithms, it is necessary to evaluate not only the level of reward attained but also the associated KL divergence. Attaining marginally higher reward at the cost of substantially increased KL divergence may not constitute an overall improvement. Figure~\ref{fig:frontier-tldr-main} displays the trade-off curve between reward and KL divergence (reward-KL frontier) across several algorithms within the sentiment task. For each algorithm, multiple runs are conducted, each utilizing a different setting for the policy’s conservativeness hyperparameter (target KL $\in\{3,6,9,12\}$ for PPO, $\beta \in \{0.05,0.1,1,5\}$ for DPO, $\alpha\in\{0.05,0.1,0.5,1\}$ for unlikelihood, and various random seeds for preferred-FT), resulting in 22 runs altogether. After every 100 training steps up to convergence, each trained policy is evaluated on a test prompt set, computing both the mean reward according to the true reward function and the average sequence-level KL divergence\footnote{That is, aggregating the KL-divergences calculated at each timestep.} from the reference policy $\text{KL}\left(\pi\mid \mid \pi_\text{ref}\right)$. Our results demonstrate that DPO yields a significantly more favorable frontier, achieving higher rewards while maintaining lower KL values than the other methods. This outcome is especially significant for several reasons. Firstly, both DPO and PPO target the same optimization objective; however, DPO shows marked improvements in efficiency, as the DPO reward/KL frontier consistently outperforms that of PPO. Secondly, DPO surpasses PPO in terms of the reward-KL trade-off frontier \emph{even when PPO is granted access to the actual ground truth rewards} (PPO-GT).

\subsection{Can DPO scale to real preference datasets?}
\label{sec:dpo-real-datasets}
We proceed to assess how well DPO performs when fine-tuned on real-world preference datasets in tasks such as summarization and single-turn dialogue. In the context of summarization, it has been shown that commonly used automatic metrics like ROUGE often do not align well with human preferences~\citep{stiennon2022learning}. Earlier studies demonstrate that leveraging human feedback through PPO for fine-tuning language models can yield more satisfactory summaries. To compare various approaches, we generate completions using the test partition of the TL;DR summarization dataset, then calculate each method's average win rate versus reference completions from the test data. Completion outputs from every approach are generated across temperature settings between 0.0 and 1.0, with their respective win rates illustrated in Figure~\ref{fig:frontier-tldr-main} (right). The models, including DPO, PPO, and Preferred-FT, are all fine-tuned starting from the same GPT-J SFT checkpoint\footnote{\url{https://huggingface.co/CarperAI/openai_summarize_tldr_sft}}. Our experiments reveal that DPO reaches a win rate of about 61\% at temperature 0.0, outperforming PPO, which achieves approximately 57\% at its optimal temperature setting. Additionally, DPO surpasses the maximum win rate attained by the best-of-$N$ baseline. It should be highlighted that we performed minimal tuning of DPO’s $\beta$ hyperparameter, suggesting the reported outcomes might underrepresent DPO’s capabilities. Furthermore, DPO displays greater robustness to changes in sampling temperature compared to PPO, whose effectiveness can decline to that of the original GPT-J model as temperature increases. Notably, Preferred-FT fails to provide substantial gains over the SFT baseline. Finally, as detailed in Section~\ref{sec:human-judgments}, we directly pit DPO and PPO against each other in human assessments: when sampling at temperature 0.25, DPO completions are favored 58\% of the time over those from PPO at the same temperature.

For single-turn dialogue evaluation, we assess various approaches using the subset of the Anthropic HH dataset's test split \citep{bai2022training}, specifically where there is a single human-assistant exchange. To compute win rates for each technique, GPT-4 comparisons reference the test set's preferred completions. In the absence of a standard SFT baseline for this benchmark, we initiate training from a pre-trained Pythia-2.8B model. We first apply Preferred-FT, fitting a reference model on selected completions to ensure outputs remain in-distribution, followed by DPO training. Our evaluation includes a comparison with the top completion out of 128 generated via Preferred-FT—the Best of $N$ method, where $N$ plateaus at 128 for this scenario (see Appendix Figure~\ref{fig:best-of-n})—and a 2-shot prompt variant of the Pythia-2.8B base model. Across all optimal temperature settings per approach, DPO matches or outperforms the alternatives. We additionally evaluate an RLHF model fine-tuned with PPO on the Anthropic HH dataset\footnote{\url{https://huggingface.co/reciprocate/ppo_hh_pythia-6B}}, utilizing code from a reputable repository\footnote{\url{https://github.com/CarperAI/trlx/tree/main/examples/hh}}; however, we are unable to identify prompt or temperature parameters that enable this model to outperform the base Pythia-2.8B. Drawing from our TL;DR results and recognizing both Best of 128 and PPO share the same reward optimization objective, we treat the former as an approximate stand-in for PPO performance. In summary, among the methods assessed, DPO uniquely offers both computational efficiency and improvement over the Anthropic HH dataset's preferred completions, achieving results on par with or superior to the resource-intensive Best of 128 baseline. Figure~\ref{fig:dialogue-main} further illustrates DPO’s rapid convergence to peak performance.

\subsection{Generalization to a new input distribution}

To investigate how PPO and DPO perform under distributional shift, we assess the policies learned from the Reddit TL;DR summarization task on an out-of-distribution benchmark: the test portion of the CNN/DailyMail dataset \citep{nallapati-etal-2016-abstractive}, which consists of news articles. We utilize the optimal sampling temperatures identified for TL;DR (0 and 0.25) when generating outputs. Table~\ref{tab:ood} presents the outcomes. The win rate is computed by having GPT-4 evaluate system summaries versus reference summaries, using the identical GPT-4 (C) prompt as in the original Reddit TL;DR assessment, with the phrase ``forum post'' replaced by ``news article.'' DPO demonstrates a clear advantage over PPO on this novel distribution. These results preliminarily indicate that DPO can generalize comparably to PPO, despite DPO not utilizing the extra unlabeled prompts from Reddit TL;DR that PPO accesses.

\subsection{Validating GPT-4 judgments with human judgments}
\label{sec:human-judgments}

We carry out a human evaluation to assess the dependability of GPT-4's scoring, employing outputs from the TL;DR summarization study and two distinct GPT-4 prompts. The \textbf{GPT-4 (S)} (simple) prompt requests the selection of the summary that better captures the main information in a post. Conversely, the \textbf{GPT-4 (C)} (concise) prompt asks for the more concise summary as well; this distinction matters since GPT-4 has a tendency to favor lengthier and more repetitive summaries under the \textbf{GPT-4 (S)} prompt compared to human evaluators. Complete prompt wordings are in Appendix~\ref{app:prompts}. We examine three methods—a top performer (DPO, temp. 0.25), a lower performer (PPO, temp. 1.0), and an intermediate baseline (SFT, temp. 0.25)—each compared against the PPO policy with its best-performing greedy sampling temperature. The comparison set was chosen to reflect a wide range of summary qualities; for DPO and PPO-1 comparisons, multiple independent human ratings were collected due to limited rater availability. Our analysis shows that both prompts lead GPT-4 to match human judgments at approximately the same rate as human inter-rater agreement, implying that GPT-4 can serve as an effective stand-in for human assessment (although our human ratings are limited in scale). Generally, the \textbf{GPT-4 (C)} prompt yields win rates that are more in line with human preferences; as such, we use this prompt for reporting main outcomes in Section~\ref{sec:dpo-real-datasets}. Further specifics about the human study design, including the rating web interface and the volunteer pool, can be found in Appendix~\ref{app:human-study}.

\section{Discussion}
Preference-based learning offers a robust and scalable approach for developing language models that are both capable and aligned with user intent. In this work, we presented DPO, a straightforward methodology that enables training language models from preferences, completely avoiding the use of reinforcement learning. Rather than adapting the preference learning task to fit within a conventional RL framework to leverage existing RL algorithms, DPO establishes a correspondence between language model policies and reward functions, allowing the model to adhere to human preferences \textit{directly} using a simple cross-entropy objective, thereby dispensing with both reinforcement learning and any loss in generality. DPO matches or exceeds the performance of traditional RLHF methods, including those employing PPO, all while requiring minimal hyperparameter adjustment. As a result, DPO significantly lowers the barrier to training additional language models guided by human preferences.

\textbf{Limitations \& Future Work.} The findings presented here open up several avenues for further research. An outstanding question is how well the DPO-trained policy performs on data outside the training distribution compared to models trained with explicit reward signals. Preliminary evidence indicates that generalization properties of DPO policies are comparable to those learned through PPO, but a more exhaustive investigation is warranted. For instance, it is worth exploring whether self-labeling using the DPO policy can also efficiently leverage unlabeled prompt data. Another area of interest concerns the manifestation of reward over-optimization in direct preference optimization, and whether the modest performance drop observed in Figure~\ref{fig:dialogue-main}-right is symptomatic of this phenomenon. Furthermore, while our experiments are limited to models containing up to 6B parameters, scaling DPO to much larger, state-of-the-art architectures remains a compelling prospect. As for evaluation procedures, our observations indicate that GPT-4's win rates are sensitive to prompt phrasing, suggesting future work should refine methods for obtaining reliable automated judgments. Finally, DPO's potential extends beyond the domain of language models trained with human preferences, offering promising opportunities for training generative models across diverse modalities.